\chapter*{Résumé}
\addcontentsline{toc}{chapter}{Résumé}

L'essor des formations à distance soulève une préoccupation croissante quant à l'honnêteté des évaluations, et c'est précisément ce problème que ce projet cherche à résoudre. Certif-fun est une plateforme qui réunit dans un seul outil la gestion des certifications, la surveillance automatisée des examens et l'accompagnement pédagogique. L'architecture repose sur des microservices distribués combinant Java Spring Boot, React et FastAPI. Pour la surveillance, des modèles de vision par ordinateur détectent en temps réel les regards suspects et les objets non autorisés, tandis qu'une vérification d'identité par QR code et reconnaissance faciale sécurise l'accès aux examens. La génération automatique de quiz s'appuie sur une IA locale via Ollama. Les tests réalisés montrent que le système détecte correctement les comportements suspects, les absences et les bruits continus, avec un temps de réponse satisfaisant dans des conditions réelles d'examen. L'interface a été jugée fluide et peu intrusive par les utilisateurs testés. Certif-fun démontre qu'il est possible de concilier rigueur sécuritaire et confort d'utilisation dans un seul outil, ouvrant la voie à des certifications en ligne dont la crédibilité peut rivaliser avec les examens en présentiel.

\bigskip
\noindent \textbf{Mots-clés :} Certification en ligne, Surveillance automatisée (Proctoring), Intégrité des évaluations, Intelligence Artificielle, Microservices, Vision par ordinateur.
